% Generated from chapter-07.md - do not edit.

\chapter{Choose To Do Something}

A better thought can change how we see a situation, but sometimes we also need to take action to help ourselves move forward. When we feel emotionally or physically stuck, it can help to think about the different choices available to us---what we can control, change, compensate for, create, or simply recognize and congratulate ourselves for doing.

\emph{Embracing Control, Change, and Growth}

\section{Action Builds Self-Esteem}

There's something very strengthening about taking action---especially small, intentional action.

One idea I like to remember is this:

\textbf{The Five C's of Choice:}

\begin{itemize}
\item \textbf{Control}
\item \textbf{Change}
\item \textbf{Compensate}
\item \textbf{Create}
\item \textbf{Congratulate}
\end{itemize}

\emph{(A nice, manageable list . . .  even before coffee.)}

Life is full of uncertainties.

We can't control our age, the way our government is run, accidents, health challenges, or the weather.

\emph{(And the weather, in particular, does not take suggestions.)}

But what we \emph{can} control is how we respond---and the choices we make every day that shape our happiness and well-being.

\section{What Can You Control?}%
\index{control}\nobreak%

While we can't control everything, we do have power over many aspects of our lives.

We manage our daily choices:

\begin{itemize}
\item where we work
\item how we care for our health
\item the friends we surround ourselves with
\item how much sleep we get
\item the food we eat
\item and even how we speak to ourselves
\end{itemize}

Instead of feeling down about not being the most loving person, we can take action---do something kind for a friend, a spouse, or a neighbor.

Rather than feeling overwhelmed by a messy house, we can clean just one room.

\emph{(Or even one drawer. Progress is progress.)}

Small, intentional steps can make a world of difference.

For example, I can choose:

\begin{itemize}
\item my sleep schedule
\item my exercise routine
\item the people I spend time with
\item and the tone I use with myself
\end{itemize}

I can choose acts of self-care---like swimming, walking, eating nutritious foods, or turning off the TV to read something uplifting.

\emph{(The TV is often surprised by this decision.)}

And while I can't control who else will be in the swimming pool---or if the power goes out---I can always adjust my mindset and look for something to appreciate.

\section{The Power to Change}%
\index{power to change}\nobreak%

Change is always within reach.

We can improve our health habits, adjust our sleep schedule, eat better, find a more fulfilling job, or make time for friends who uplift us.

A woman once shared in an evaluation:

``I learned that it wasn't my fault---and I could change.''

That realization is powerful.

One thing that still amazes me is this:

\textbf{I can change my beliefs.}

And when I change my beliefs, I can change my behavior.

During the class I took with Jackie Kelm, I discovered that I could change the belief that food temptation was too powerful for me.\index{Jackie Kelm}

Now, I remind myself each day:

\textbf{I have food mastery.}\index{food mastery}

\emph{(Sometimes daily reminders are required. That's perfectly fine.)}

We can also change:

\begin{itemize}
\item the way we speak to ourselves
\item the environment we live in
\item the habits we practice
\end{itemize}

We can move to a new apartment, find a better doctor, ask for help rearranging furniture, or simply ask a friend for support.

Even small situations can shift.

If we're stuck in traffic, we can turn frustration into a mini-meditation.

\emph{(Not always easy---but available.)}

\section{Pause for a Moment}

What is one small thing you \emph{can} control right now?

A thought?\\
A word you say?\\
A tiny action?

You don't have to change everything.

Just one small choice is enough.

\section{Compensate for Life's Challenges}%
\index{challenges}\nobreak%

Loss and disappointment are part of life.

But we can \textbf{compensate}---we can gently fill the gaps with care and support.

If we've lost a loved one, we can lean into friendship and connection.

If work isn't fulfilling, we can create meaning through hobbies or relationships.

If a difficult interaction leaves us unsettled, we can respond with self-compassion.

By replacing negative messages with kinder ones, we take back our balance.

Speak to yourself the way you would speak to someone you love.

\emph{(You deserve that same tone.)}

\section{Create Joy in Everyday Life}%
\index{joy everyday}\nobreak%

We have more influence over our days than we sometimes realize.

We can add or subtract activities to make life more enjoyable:

\begin{itemize}
\item meeting friends
\item listening to music
\item watching something inspiring
\item joining a group
\item or simply sitting quietly with a cup of tea
\end{itemize}

Our only real limitation is our energy---and we can choose how to invest it.

\emph{(Energy is precious. Spend it wisely . . .  and kindly.)}

\section{Congratulate Yourself}%
\index{congratulate yourself}\nobreak%

This is a crucial step---and one we often forget.

A friend of mine once asked:

\textbf{``Do you cheer for yourself when you do something good?''}

What a wonderful idea.

Recognizing and celebrating our own efforts is one of the kindest things we can do for ourselves.

\begin{itemize}
\item Acknowledge your progress
\item Celebrate small victories
\item Count your blessings
\end{itemize}

Reward yourself after a tough moment---like getting through a dental appointment . . . 

or surviving the DMV.

\emph{(Both worthy of recognition.)}

\section{A Gentle Closing Reflection}

Life is not always easy.

But through small, thoughtful choices, we can shape it into something more manageable---and more joyful.

So take a moment right now:

\textbf{What's one thing you can celebrate about yourself today?}

Even something small counts.

\emph{(Especially something small.)}

\section{A Simple Reminder}

You can be in action to raise your self-esteem.

Remember the Five C's:

\textbf{Control. Change. Compensate. Create. Congratulate.}

And underneath all of it, a quiet intention:

\textbf{I choose.}

I choose how I respond.\\
I choose how I speak to myself.\\
I choose the kind of person I want to be.

And sometimes, that choice is very simple:

\textbf{Happy.}
